\begin{proposition}[the marked genus-three component cover]\label{push:src:B.3.2}\label{push:B32}
Let $B^\circ=\mathbb P^1\setminus\{0,1,\infty\}$ and use the
clockwise deck-identification convention of \ref{push:src:B.3.1}.
The peripheral rotations $4,-3,-1$ in $\mathbb Z/12$ define a
connected degree-twelve cover of $B^\circ$. Its compact
normalization is the smooth genus-three curve obtained from
\[
 y^{12}=\frac{t^4}{(t-1)^3}.
\]
The ramification indices over $0,1,\infty$ are $3,4,12$,
respectively. The peripheral relation determines the last rotation
from the first two, so this is the marked cyclic component cover
used in the partial filling of \ref{push:src:B.4.19}.
\end{proposition}
\Uses{\cref{push:src:B.3.1}; the explicit function-field equation and its local normalizations}
\begin{proof}
The presentation $g_1g_2g_0=1$ sends the third peripheral generator
to $-4+3=-1$. The first two images generate $\mathbb Z/12$, so
the resulting action on the twelve sheets is transitive. This
constructs the connected cover by its monodromy representation.

The rational function $t^4/(t-1)^3$ has orders $4,-3,-1$ at the
three punctures. Its twelfth-root equation has the same marked
deck representation. The covering-space identification is the map
which sends a chosen twelfth root at one base point to its chosen
sheet and continues both along the same paths. Its equality of
monodromy makes that map single-valued; its local analytic root
coordinates make it holomorphic. Normalization extends this
identification to the compact covers.

Locally $y^{12}=z^a u(z)$ has $\gcd(12,a)$ branches of ramification
index $12/\gcd(12,a)$, after absorbing the holomorphic unit in a
local twelfth root. Thus the three contributions to ramification
are $4(3-1)$, $3(4-1)$, and $12-1$. Their sum is $28$.
Riemann--Hurwitz gives $2-2g=2\cdot12-28=-4$, so $g=3$.
Reversing the deck convention is realized by $y\mapsto y^{-1}$,
with the reciprocal defining equation; the compact curve and the
ramification calculation are unchanged.
\end{proof}
